\chapter{The Chinese remainder theorem and Fermat--Euler}\label{ch:crt}
\index{Chinese remainder theorem}\index{Fermat's little theorem}

Chapter~6 solved one congruence at a time. This chapter joins them ---
the Chinese remainder theorem says a system of congruences with coprime
moduli always has exactly one solution class mod the product --- and then
proves the two big structure theorems of prime moduli: Fermat's little
theorem and Euler's totient theorem. These are the workhorses behind
almost every ``compute a huge power mod $m$'' problem in the archive.
After this chapter you will be able to solve CRT systems by construction,
reduce exponents by the right period ($\varphi$, not $m$), and recognise
when a congruence problem is really asking for an order.

\section{The Chinese remainder theorem}\label{sec:crt}
\index{CRT|see{Chinese remainder theorem}}

\begin{theorem}[CRT]\label{thm:crt}
If $m_1, \dots, m_k$ are pairwise coprime, the system $x \equiv r_i
\pmod{m_i}$ has exactly one solution modulo $M = m_1 m_2 \cdots m_k$.
\end{theorem}

Existence \emph{is} the construction, and the construction is the two-
modulus case repeated: to merge $x \equiv r_1 \pmod{m_1}$ with
$x \equiv r_2 \pmod{m_2}$, write $x = r_1 + m_1 t$ and demand
$m_1 t \equiv r_2 - r_1 \pmod{m_2}$ --- solvable because
$\gcd(m_1, m_2) = 1$ gives $m_1$ an inverse (Theorem~\ref{thm:cancel}).
Uniqueness: two solutions differ by a multiple of every $m_i$, and
pairwise coprimality makes that a multiple of $M$.

\begin{example}\label{ex:crt}
Solve $x \equiv 2 \pmod 3$, $x \equiv 3 \pmod 5$, $x \equiv 2 \pmod 7$.
\par\smallskip
\textbf{Strategy.} Satisfy one modulus, then \emph{add multiples} to
satisfy the next: parametrise the current solution as $x = r + mt$ and
solve the linear congruence for $t$ --- Chapter~6's entire toolbox inside
Chapter~7's loop.
\par\smallskip
\begin{solution}
$x \equiv 2 \pmod 3$: write $x = 2 + 3t$. Mod $5$: $2 + 3t \equiv 3$,
$3t \equiv 1$, $t \equiv 2$ --- so $x \equiv 8 \pmod{15}$. Fold again:
$x = 8 + 15s$; mod $7$: $8 + 15s \equiv 2$, i.e.\ $15s \equiv -6 \equiv 1 \pmod 7$ ---
and $15 \equiv 1 \pmod 7$, so $s \equiv 1$.
Hence $x \equiv 23 \pmod{105}$ --- the unique class, verified in the
listing below ($23 \bmod 3 = 2$, $\bmod 5 = 3$, $\bmod 7 = 2$).
\end{solution}
\end{example}

Figure~\ref{fig:rulers} draws the two-modulus view: residues on two
different rulers, and the solution is the tick where both land.

\begin{figure}[ht]\centering
% Figure: CRT as two sliding rulers --- residues mod 3 and mod 5.
% Source for Figure~\ref{fig:rulers}. TikZ only; regenerable by compiling main.tex.
\begin{tikzpicture}[
  tick/.style={draw=rulegray},
  hit/.style={circle,fill=accent,inner sep=1.6pt},
  lab/.style={font=\footnotesize\sffamily,text=slate},
]
% top ruler: positions 0..14 labeled mod 3
\node[lab,anchor=east] at (-0.3,0) {mod 3:};
\foreach \i in {0,...,14}
{
  \pgfmathtruncatemacro{\v}{mod(\i,3)}
  \draw[tick] (\i*0.55,0) -- (\i*0.55,0.18);
  \node[lab] at (\i*0.55,0.42) {\v};
  \ifnum\v=2 \node[hit] at (\i*0.55,0) {}; \fi
}
% bottom ruler: positions mod 5
\node[lab,anchor=east] at (-0.3,-1.0) {mod 5:};
\foreach \i in {0,...,14}
{
  \pgfmathtruncatemacro{\v}{mod(\i,5)}
  \draw[tick] (\i*0.55,-1.0) -- (\i*0.55,-0.82);
  \node[lab] at (\i*0.55,-0.58) {\v};
  \ifnum\v=3 \node[hit] at (\i*0.55,-1.0) {}; \fi
}
% alignment markers at 8 and 23->position 8
\draw[accent,dashed] (8*0.55,0.15) -- (8*0.55,-1.15);
\node[lab] at (8*0.55,-1.55) {$x = 8$: hits both conditions};
\node[lab] at (8*0.55,-1.95) {(and again at $x = 8 + 15$, $8 + 30$, \dots)};
\end{tikzpicture}

\caption{CRT as sliding rulers: $x \equiv 2 \pmod 3$ (top dots) and
$x \equiv 3 \pmod 5$ (bottom dots) meet at $x = 8$, and again every $15$
--- the period is the \emph{product}, which is why coprimality matters.}
\label{fig:rulers}
\end{figure}

\begin{checkpoint}[Predict before you run]
Listing~\ref{lst:crt} solves the system two ways --- brute scan and the
folding construction. Predict the intermediate fold $x \equiv ? \pmod{15}$
before reading the output.
\end{checkpoint}

\begin{lstlisting}[language=Python,caption={\texttt{crt\_solver.py}: fold
two congruences at a time via B\'ezout (full script in
\texttt{code/crt\_solver.py}).},label={lst:crt}]
def crt_pair(r1, m1, r2, m2):
    g, s, _ = egcd(m1, m2)
    t = (s * (r2 - r1) // g) % (m2 // g)
    return (r1 + m1 * t) % (m1 * m2 // g), m1 * m2 // g
\end{lstlisting}
Real output:
\begin{lstlisting}[numbers=none]
system: x == 2 (mod 3), x == 3 (mod 5), x == 2 (mod 7)
brute force scan: [23]
  folded: x == 8 (mod 15)
  folded: x == 23 (mod 105)
constructive: x == 23 (mod 105)
unique mod 105: 23; check [2, 3, 2]
\end{lstlisting}

\begin{exercise}[easy]\label{exr:crt2}
Solve $x \equiv 1 \pmod 4$, $x \equiv 4 \pmod 9$ by the folding method,
and verify against \texttt{crt\_solver.py}.
\end{exercise}

\section{Three ways to hold the theorem}\label{sec:lenses}
\index{CRT!three views}

Chen's handout presents CRT as three lenses on the same fact, and each
lens wins a different fight~\cite{chen-crt}:
\begin{itemize}
\item \textbf{Construction} (what we just did): given target residues,
\emph{build} an $x$ with them. Use when the problem asks ``find an
integer such that \dots'' --- the theorem licenses your assembly.
\item \textbf{Lifting}: if $x \equiv r$ holds mod every part, it holds
mod the product. $x \equiv 1$ mod $2,3,5$ means $x \equiv 1 \pmod{30}$ ---
surprising congruences often decompose this way.
\item \textbf{Destruction}: a congruence mod $M$ can be \emph{split} into
its prime-power parts, which are easier to reason about one at a time.
``$a^2 \equiv 1 \pmod{24}$'' is really ``mod $8$ and mod $3$
simultaneously'' --- and each half is nearly free.
\end{itemize}
The skill is recognising which direction the problem is already facing:
systems want construction, big moduli want destruction, and ``show it
holds for all $n$'' often wants lifting.

\section{Systems in practice}\label{sec:crtpractice}
\index{CRT!non-coprime moduli}

The theorem's hypothesis is load-bearing. With $\gcd(m_1, m_2) = d > 1$,
a solution exists \emph{iff} $r_1 \equiv r_2 \pmod d$ --- the consistency
check --- and then the class is unique only mod
$\mathrm{lcm}(m_1, m_2)$. $x \equiv 1 \pmod 4$, $x \equiv 3 \pmod 6$:
consistent ($1 \equiv 3 \pmod 2$), solutions $x \equiv 9 \pmod{12}$ ($9 \bmod 4 = 1$, $9 \bmod 6 = 3$). But
$x \equiv 1 \pmod 4$, $x \equiv 4 \pmod 6$ is inconsistent: no solution,
and no amount of B\'ezout saves it. The gate is cheap --- always run it
before constructing.

\begin{exercise}[easy]\label{exr:crtgate}
Decide, without solving: is $x \equiv 2 \pmod 6$, $x \equiv 5 \pmod 9$
consistent? If yes, solve it; if no, say what the gate detected.
\end{exercise}

\section{Fermat's little theorem, two ways}\label{sec:flt}
\index{Fermat's little theorem}

\begin{theorem}[Fermat's little theorem]\label{thm:flt}
For prime $p$ and $p \nmid a$: $a^{p-1} \equiv 1 \pmod p$. Equivalently
$a^p \equiv a \pmod p$ for all $a$.
\end{theorem}

\begin{example}\label{ex:flt}
Prove it by the multiplication-permutation argument.
\par\smallskip
\textbf{Strategy.} Multiply the nonzero residues $1, 2, \dots, p-1$ by
$a$: because $a$ is coprime to $p$, that map is a \emph{bijection} ---
the same set reshuffled (Figure~\ref{fig:perm}). Two lists with the same
members have the same product; the rest is cancellation.
\par\smallskip
\begin{solution}
$\{a \cdot 1, a \cdot 2, \dots, a\cdot(p-1)\} \pmod p$ are the nonzero
residues in some order (no two collide: $ai \equiv aj$ would cancel $a$
--- coprimality again). Hence
$a \cdot 2a \cdots (p-1)a \equiv 1 \cdot 2 \cdots (p-1) \pmod p$, i.e.\
$a^{p-1} (p-1)! \equiv (p-1)!$. Since $\gcd((p-1)!, p) = 1$, cancel:
$a^{p-1} \equiv 1 \pmod p$.
\end{solution}
\end{example}

\begin{figure}[ht]\centering
% Figure: multiplication by 3 permutes nonzero residues mod 7.
% Source for Figure~\ref{fig:perm}. TikZ only; regenerable by compiling main.tex.
\begin{tikzpicture}[
  nd/.style={circle,draw,minimum size=6.5mm,font=\small},
  ed/.style={->,accent,thick},
  lab/.style={font=\footnotesize\sffamily,text=slate},
]
\foreach \i in {1,...,6}
  \node[nd] (n\i) at ({90 - (\i-1)*60}:2.0) {\i};
% times 3 mod 7: 1->3, 3->2, 2->6, 6->4, 4->5, 5->1
\draw[ed,bend left=14] (n1) to (n3);
\draw[ed,bend left=14] (n3) to (n2);
\draw[ed,bend left=14] (n2) to (n6);
\draw[ed,bend left=14] (n6) to (n4);
\draw[ed,bend left=14] (n4) to (n5);
\draw[ed,bend left=14] (n5) to (n1);
\node[lab] at (3.0,0.4) {$\times 3 \bmod 7$: one 6-cycle};
\node[lab] at (3.0,0.0) {a bijection $\Rightarrow$};
\node[lab] at (3.0,-0.4) {products match:};
\node[lab] at (3.0,-0.8) {$\prod (3r) \equiv \prod r$};
\node[lab] at (0,-2.9) {that product equality is Fermat's little theorem};
\end{tikzpicture}

\caption{$\times 3$ mod $7$ is a bijection --- here a single 6-cycle.
Notice the proof hiding in the picture: the products of all arrows'
endpoints and origins are the same set, so $3^6 \cdot 6! \equiv 6!$.}
\label{fig:perm}
\end{figure}

\begin{example}[FLT by counting]
Prove $a^p \equiv a \pmod p$ by a combinatorial argument.
\par\smallskip
\textbf{Strategy.} Count $a^p$ strings over $a$ colours, remove the $a$
monochrome ones, and partition the rest into rotation classes ---
primality is what forces every class to have full size $p$.
\par\smallskip
\begin{solution}
Of the $a^p$ strings on an alphabet of size $a$, exactly $a$ are
monochrome. Every other string has all $p$ of its rotations distinct
(a smaller orbit would give a nontrivial divisor of $p$ --- see
Figure~\ref{fig:necklace}), so the remaining $a^p - a$ strings split
into classes of size $p$. Hence $p \mid a^p - a$.
\end{solution}
\end{example}

\begin{figure}[ht]\centering
% Figure: necklaces for the FLT counting proof --- p beads, a colours.
% Source for Figure~\ref{fig:necklace}. TikZ only; regenerable by compiling main.tex.
\begin{tikzpicture}[
  bead/.style={circle,draw=ink,minimum size=5mm,font=\scriptsize},
  mono/.style={bead,fill=accentsoft},
  lab/.style={font=\footnotesize\sffamily,text=slate},
]
% 5-bead necklace, pattern A B A A C
\foreach \i/\c in {0/A,1/B,2/A,3/A,4/C}
  \node[bead] at ({90+\i*72}:1.1) {\c};
\node[lab] at (0,-1.9) {one necklace, $p = 5$ beads, $a$ colour names};
% three monochrome necklaces (small rings at x = 3.0, 4.3, 5.6)
\foreach \k in {0,1,2}
{
  \foreach \i in {0,...,4}
  {
    \pgfmathsetmacro{\px}{3.0 + \k*1.3 + 0.42*cos(90+\i*72)}
    \pgfmathsetmacro{\py}{0.42*sin(90+\i*72)}
    \node[mono] at (\px,\py) {};
  }
  \node[lab] at (3.0+\k*1.3,-1.0) {\ifnum\k=0 all-A\else\ifnum\k=1 all-B\else all-C\fi\fi};
}
\node[lab] at (4.3,-1.9) {$a$ monochrome necklaces};
\node[lab] at (0,-2.7) {the other $a^p - a$ strings sit in rotation-classes of size $p$};
\node[lab] at (0,-3.1) {($p$ prime: an orbit smaller than $p$ would repeat a proper divisor)};
\node[lab] at (0,-3.5) {$\Rightarrow\; a^p - a = p\cdot(\text{classes}) \equiv 0 \pmod p$};
\end{tikzpicture}

\caption{The necklace proof of $a^p \equiv a \pmod p$. Notice why $p$
must be prime: only then does every non-monochrome string rotate into
\emph{exactly} $p$ distinct strings --- a composite modulus would allow
smaller orbits and break the count.}
\label{fig:necklace}
\end{figure}

\begin{exercise}[medium]\label{exr:flt}
Compute $2^{100} \bmod 13$ and $7^{222} \bmod{11}$ by Fermat reduction;
then verify with \texttt{powers\_cycle.py}.
\end{exercise}

\section{Euler's totient and Euler's theorem}\label{sec:euler}
\index{totient}\index{Euler's theorem}

\begin{definition}
$\varphi(n)$ counts the integers in $\{1, \dots, n\}$ coprime to $n$.
For $n = \prod p_i^{e_i}$: $\varphi(n) = n \prod (1 - 1/p_i) =
\prod p_i^{e_i-1}(p_i - 1)$.
\end{definition}

\begin{theorem}[Euler]\label{thm:euler}
If $\gcd(a, n) = 1$: $a^{\varphi(n)} \equiv 1 \pmod n$.
\end{theorem}

The proof is Fermat's verbatim, with ``nonzero residues'' replaced by
``units'' --- the $\varphi(n)$ residues coprime to $n$ form a group under
multiplication, and multiplying by a unit permutes them exactly as in
Figure~\ref{fig:perm}. Fermat is the special case $\varphi(p) = p - 1$.
Figure~\ref{fig:totient} charts $\varphi$ to $36$: the primes are the
spikes where $n-1$ is achieved.

\begin{figure}[ht]\centering
% Figure: totient phi(n) for n = 1..36 --- primes spike, products echo.
% Source for Figure~\ref{fig:totient}. TikZ only; regenerable by compiling main.tex.
\begin{tikzpicture}[
  bar/.style={fill=accentsoft,draw=rulegray},
  pbar/.style={fill=accent,draw=accent},
  lab/.style={font=\scriptsize\sffamily,text=slate},
]
\foreach \n/\f in {1/1,2/1,3/2,4/2,5/4,6/2,7/6,8/4,9/6,10/4,11/10,12/4,
                   13/12,14/6,15/8,16/8,17/16,18/6,19/18,20/8,21/12,22/10,
                   23/22,24/8,25/20,26/12,27/18,28/12,29/28,30/8,31/30,
                   32/16,33/20,34/16,35/24,36/12}
{
  \pgfmathtruncatemacro{\isp}{\f == \n-1 ? 1 : 0}
  \ifnum\isp=1
    \draw[pbar] (\n*0.24,0) rectangle (\n*0.24+0.18,\f*0.06);
  \else
    \draw[bar] (\n*0.24,0) rectangle (\n*0.24+0.18,\f*0.06);
  \fi
}
\draw[draw=rulegray] (0.2,0) -- (8.9,0);
\node[lab,anchor=north] at (1.1,-0.1) {5};
\node[lab,anchor=north] at (2.3,-0.1) {10};
\node[lab,anchor=north] at (3.5,-0.1) {15};
\node[lab,anchor=north] at (4.7,-0.1) {20};
\node[lab,anchor=north] at (5.9,-0.1) {25};
\node[lab,anchor=north] at (7.1,-0.1) {30};
\node[lab,anchor=north] at (8.3,-0.1) {35};
\node[lab] at (4.5,2.2) {dark bars: $\varphi(p) = p - 1$ exactly at primes};
\node[lab] at (4.5,-0.85) {notice $\varphi(p) = p-1$ is the record --- composite $n$ wastes factors};
\end{tikzpicture}

\caption{$\varphi(n)$ for $n \le 36$. Notice the records: $\varphi(p) =
p-1$ at primes (dark bars), and values like $\varphi(30) = 8$ show the
product discount $30(1-\tfrac12)(1-\tfrac13)(1-\tfrac15)$ at work.}
\label{fig:totient}
\end{figure}

\begin{example}\label{ex:last3}
Find the last three digits of $3^{400}$.
\par\smallskip
\textbf{Strategy.} Last three digits is mod $1000$; Euler reduces the
exponent mod $\varphi(1000)$. Compute $\varphi(1000) = 1000(1 -
\tfrac12)(1 - \tfrac15) = 400$ --- the exponent \emph{is} the order's
multiple.
\par\smallskip
\begin{solution}
$\gcd(3, 1000) = 1$, so $3^{400} = 3^{\varphi(1000)} \equiv 1
\pmod{1000}$: the last three digits are \textbf{001}.
\end{solution}
\end{example}

\begin{exercise}[medium]\label{exr:euler}
Find the last three digits of $7^{1000}$ and of $3^{401}$ (reduce each
exponent mod the period from WE~\ref{ex:last3} first; mind the
remainder this time).
\end{exercise}

\section{Orders modulo $p$}\label{sec:orders}
\index{order!multiplicative}

Euler gives \emph{an} exponent that returns $1$; the \emph{order} of $a$
mod $m$ is the \emph{smallest} positive one, written $\operatorname{ord}_m(a)$.
Two facts carry the theory: $\operatorname{ord}$ always divides
$\varphi(m)$ (long division of the exponent), and $a^k \equiv 1$ iff the
order divides $k$. Chapter~6's power cycles were orders in plain clothes
--- $7$ has order $4$ mod $100$, and $4 \mid \varphi(100) = 40$. A unit
whose order is $\varphi(m)$ is a \emph{primitive root}; mod primes one
always exists (proved in the olympiad track~\cite{zhao}).

\begin{exercise}[medium]\label{exr:ord}
Find $\operatorname{ord}_{11}(2)$ and $\operatorname{ord}_{11}(10)$;
both divide $10$ --- which hits the bound, and what does that mean?
\end{exercise}

\section{Wilson's theorem (enrichment)}\label{sec:wilson}
\index{Wilson's theorem}

\begin{theorem}
$(p-1)! \equiv -1 \pmod p$ for prime $p$.
\end{theorem}

The proof pairs each residue with its inverse: mod $p$, the only
self-inverse units are $\pm 1$ ($x^2 \equiv 1$ factors as
$p \mid (x-1)(x+1)$, and primality gives no zero divisors). Every other
unit cancels in pairs, leaving $1 \cdot (-1)$. The converse is also true
--- the congruence \emph{characterises} primes --- but useless as a
primality test ($p!$ is huge); its real job is inside other proofs.

\section{Applications: last digits again, check digits}\label{sec:apps}
\index{check digit}

Chapter~6 computed last digits by hunting periods; Euler supplies the
upper bound on the hunt. And the same machinery hides in every checkout
line: the ISBN-10 check digit $d_{10} = \sum_{i=1}^{9} i \cdot d_i \bmod
{11}$ catches \emph{any} single-digit error and any transposition of
adjacent digits --- the second claim is exactly that the weights
$i$ are distinct mod $11$.

\begin{example}\label{ex:usamo}
Show that $t^2 + t + 1$ has arbitrarily many distinct prime divisors as
$t$ ranges over the positive integers.
\par\smallskip
\textbf{Strategy.} Euclid's infinitude argument, lifted onto the
value-set: given any finite list of primes hitting the values, evaluate
$f$ at their product --- every listed prime divides the argument, so the
value is $\equiv f(0) = 1$ at each one, and every prime factor of the
value is new.
\par\smallskip
\begin{solution}
Suppose only $p_1, \dots, p_k$ ever divide values of
$f(t) = t^2 + t + 1$. Set $Q = p_1 p_2 \cdots p_k$ and evaluate
$f(Q) = Q^2 + Q + 1$. For each listed $p_i$: $Q \equiv 0 \pmod{p_i}$, so
$f(Q) \equiv f(0) = 1 \pmod{p_i}$ --- no listed prime divides $f(Q)$.
But $f(Q) > 1$ has some prime factor $q$, which is therefore a new prime
dividing a value of $f$: contradiction. The same trick works for any
nonconstant polynomial with $f(0) = \pm 1$ (USAMO 2008 P1 builds the full
construction on this idea~\cite{chen-crt}).
\end{solution}
\end{example}

\section{Chapter summary}
\begin{table}[ht]\centering
\begin{tabularx}{\linewidth}{@{}lYY@{}}
\toprule
\tablehead\thead{Tool} & \thead{Condition} & \thead{Delivers}\\
\midrule
CRT (coprime) & $m_i$ pairwise coprime & one class mod $\prod m_i$\\
CRT (general) & consistent iff $r_i \equiv r_j \pmod{\gcd}$ & class mod $\mathrm{lcm}$\\
Fermat & $p$ prime, $p \nmid a$ & $a^{p-1} \equiv 1 \pmod p$\\
Euler & $\gcd(a,n)=1$ & $a^{\varphi(n)} \equiv 1 \pmod n$\\
totient & $n = \prod p^e$ & $n\prod(1-1/p)$; never trial-count\\
order & $\operatorname{ord} \mid \varphi(m)$ & exact period, divides $\varphi$\\
last digits of $a^N$ & reduce $N$ mod $\operatorname{ord}$, not $m$ & finite loop\\
\bottomrule
\end{tabularx}
\caption{The chapter in one page: every tool is a gate plus a formula.
The second column is where the proofs actually check you.}
\label{tab:crt}
\end{table}

\section{Common mistakes}
\index{common mistakes}
\begin{itemize}
\item \textbf{CRT on non-coprime moduli.} The construction proceeds
whether or not the system is consistent --- and returns a confident wrong
answer. Gate first: $r_i \equiv r_j \pmod{\gcd(m_i, m_j)}$ for every
pair, else the system is just unsatisfiable.
\item \textbf{Fermat off-by-one.} $a^{p-1} \equiv 1$ needs $p \nmid a$
\emph{and} exponent $p-1$ --- writing $a^p \equiv 1$ is the classic slip
(the correct all-$a$ form is $a^p \equiv a$). Reduce exponents mod
$p-1$, not mod $p$.
\item \textbf{Totient by trial.} Counting coprimes to $n$ by hand is
$O(n)$ and error-prone; the product $n\prod(1-1/p)$ needs only the
factorisation. $\varphi(1000) = 400$ in one line, not 400 checks.
\end{itemize}

\section{Additional exercises}
\begin{exercise}[easy]
Solve $x \equiv 1 \pmod 2$, $x \equiv 2 \pmod 3$, $x \equiv 3 \pmod 5$
by folding; verify with \texttt{crt\_solver.py}.
\end{exercise}
\begin{exercise}[easy]
Show $n^5 \equiv n \pmod{30}$ for every $n$ --- split into the three
moduli (destruction lens) and use Fermat on each.
\end{exercise}
\begin{exercise}[easy]
Compute $\varphi(360)$ and $\varphi(97)$; list the units of $\Z/8\Z$ and
verify each squares to $1$.
\end{exercise}
\begin{exercise}[easy]
Find the last two digits of $3^{100}$ via Euler ($\varphi(100)$) and via
the cycle picture of Chapter~6 --- confirm they agree.
\end{exercise}
\begin{exercise}[medium]
Prove the totient product formula for prime powers first:
$\varphi(p^e) = p^e - p^{e-1}$ by direct count; then argue
multiplicativity from CRT (a unit mod $mn$ is a pair of units).
\end{exercise}
\begin{exercise}[medium]
Solve $x \equiv 4 \pmod 6$, $x \equiv 2 \pmod 8$ after checking
consistency; state the modulus of the answer class.
\end{exercise}
\begin{exercise}[medium]
Prove $n^{13} \equiv n \pmod{2730}$ for all $n$ --- the lcm of
$\{p-1 : p \mid 2730\}$ is the engine.
\end{exercise}
\begin{exercise}[medium]
Show that if $\operatorname{ord}_m(a) = k$, then $a^i \equiv a^j$ iff
$i \equiv j \pmod k$ --- the power cycle is exactly a cycle.
\end{exercise}
\begin{exercise}[hard]
Find all $n \le 60$ with $\varphi(n) = 12$, and prove your list is
complete (what does $p - 1 \mid 12$ force?).
\end{exercise}
\begin{exercise}[hard]
Prove that $x^2 \equiv 1 \pmod p$ has \emph{only} the solutions
$x \equiv \pm 1$ --- and that this fails mod $8$: find all four.
\end{exercise}
\begin{exercise}[hard]
Extend \texttt{crt\_solver.py} to non-coprime systems: run the
consistency gate, emit the class mod $\mathrm{lcm}$, and report
inconsistent systems explicitly. Verify on $x \equiv 4 \pmod 6$,
$x \equiv 2 \pmod 8$ and a deliberately broken system.
\end{exercise}

\chapternotes
The three-lenses framing (construction, lifting, destruction) and the
USAMO~2008~P1 connection follow Chen's CRT handout~\cite{chen-crt}. The
folding construction and non-coprime consistency gate follow 6.042J's
CRT sections~\cite{mcs}; Velleman supplied the residue-system
proofs~\cite{velleman}; Berkeley CS70's notes push on into RSA from
exactly this point~\cite{cs70}, and Zhao's handouts carry orders and
primitive roots to olympiad depth~\cite{zhao}. The necklace proof of
Fermat is the standard counting argument (cf.\ OTIS~\cite{otis}).
Contest items this chapter powers --- CRT systems (AIME~2023~I~P7,
AIME~2025~II~P2), Euler reductions (AIME~2024~I~P13), remainder cycles
(AIME~2020~I~P12), orders and $\varphi$ (AIME~2010~I~P14,
AIME~2021~I~P7), plus USAJMO~2024~P3 and USAMO~2008~P1 --- are on the
AoPS archive~\cite{aops-amc}.
